\documentclass[11pt]{article}

\usepackage[letterpaper,margin=1in]{geometry}
\usepackage[T1]{fontenc}
\usepackage[utf8]{inputenc}
\usepackage{lmodern}
\usepackage{amsmath}
\usepackage{microtype}
\usepackage{setspace}
\usepackage{enumitem}
\PassOptionsToPackage{hyphens}{url}
\usepackage[colorlinks=true,linkcolor=blue,citecolor=blue,urlcolor=blue]{hyperref}

\setstretch{1.08}
\setlength{\parindent}{1.5em}
\setlength{\parskip}{0.25em}
\setlist[itemize]{leftmargin=1.5em}
\setlist[enumerate]{leftmargin=1.5em}

\providecommand{\tightlist}{%
  \setlength{\itemsep}{0pt}\setlength{\parskip}{0pt}}

\title{Governing Trustworthy Agentic Healthcare AI:\\
A Calibration-First Framework with Distributed Ledger Accountability}
\author{Michael O. Eniolade\\
University of the Cumberlands\\
\texttt{meniolade20593@ucumberlands.edu}}
\date{May 2026}

\begin{document}
\maketitle

\noindent\textbf{Keywords:} agentic AI; healthcare governance; probability calibration; distributed ledger technology; clinical decision support; trustworthy AI.

\begin{abstract}
Agentic artificial intelligence in healthcare is advancing faster than
the governance frameworks designed to oversee it. These systems are
autonomous, goal-directed, and capable of multi-step reasoning and
sequential decision-making, introducing accountability challenges that
existing clinical AI standards were not built to address. TRIPOD+AI,
DECIDE-AI, and the FDA's AI/ML-based Software as a Medical Device
guidance evaluate static models at a point in deployment. They do not
address pipeline-level reliability, compounding errors across agent
decision steps, or the accountability vacuum created when autonomous
systems act at speeds beyond human review.

This perspective proposes a five-layer trustworthiness governance
framework for agentic healthcare AI: (1) calibration integrity and
uncertainty quantification at each decision point; (2) subgroup
fairness auditing; (3) verifiable model provenance; (4)
patient-controlled consent management for autonomous AI actions; and
(5) scored deployment readiness with ongoing monitoring. Each layer
carries measurable thresholds: Expected Calibration Error at or below
0.10, conformal prediction coverage at or above 0.90, and a subgroup
ECE gap below 0.05.

Blockchain and distributed ledger technologies (DLT) provide the
accountability infrastructure that makes trustworthiness evidence
verifiable, persistent, and independently auditable. Immutable audit
trails, smart contract-encoded constraints, and decentralized consent
registries address the epistemic accountability gap that traditional
logging cannot fill. The framework connects to emerging research on
clinical AI calibration, conformal prediction, and deployment readiness
evaluation, and is designed for empirical validation in subsequent
research.
\end{abstract}


\section*{Plain Language Summary}\label{plain-language-summary}

\textbf{What this paper is about.} Healthcare AI is going ``agentic.''
These systems make sequences of decisions without a human approving each
step. An AI agent might monitor a patient's vital signs, detect a
deterioration pattern, and automatically alert the care team or initiate
escalation. These systems are already in clinical use.

\textbf{The problem.} The rules we use to keep healthcare AI safe were
designed for tools that make one decision at a time. Those older rules
do not address what happens when AI makes many consecutive decisions,
how errors compound across steps, or who is accountable when an
autonomous system causes harm.

\textbf{What we propose.} This paper proposes five requirements for any
agentic AI system in healthcare: (1) probability predictions should be
well-calibrated at every decision step, (2) the system should work
equitably across patient groups, (3) training history and data sources
should be traceable and verifiable, (4) patients should have control
over what AI is permitted to do on their behalf, and (5) the system
should pass a structured readiness assessment before deployment, with
continuous monitoring thereafter.

\textbf{Why blockchain helps.} Blockchain creates permanent,
tamper-proof digital records. This provides the foundation to make these
five requirements enforceable and independently verifiable.
Blockchain-based audit trails, smart contracts, and decentralized
consent registries turn governance rules from policy documents into
operational constraints built into the system itself.

\textbf{Who this is for.} Clinical informaticists, health system
leaders, AI developers, and policymakers who need a practical governance
framework for deploying agentic AI responsibly in healthcare.


\section{Introduction}\label{introduction}

\subsection{The Agentic Shift in Healthcare
AI}\label{the-agentic-shift-in-healthcare-ai}

Artificial intelligence in healthcare has crossed a threshold. For more
than a decade, clinical AI research focused on a well-defined paradigm:
train a model on historical data, validate it against a held-out test
set, report discrimination metrics, and evaluate whether predictions
exceed a performance threshold. The fundamental unit was a single
inference: one patient, one input, one prediction. Clinical
decision-support systems operating in this mode function as
sophisticated calculators, responding to a clinician's query with a risk
score or classification label.

This paradigm is being supplanted by a qualitatively different model of
AI deployment. Agentic AI systems are autonomous, goal-directed software
entities capable of multi-step reasoning, tool use, planning, and
sequential action. They do not respond to a single prompt and stop. They
pursue goals, decompose tasks, invoke external systems, retrieve
information, make intermediate decisions, and produce outputs that often
trigger further automated actions. Clinical examples span multiple care
settings. Some systems autonomously monitor patient deterioration
signals across data streams and initiate escalation protocols. Others
iterate through laboratory and imaging results to refine differential
diagnoses. Medication safety agents cross-check orders against patient
records, interaction databases, and renal function values. Care
coordination systems schedule and adjust clinical workflows without a
clinician approving each step~{[}1{]}. A 2025 systematic review of AI
agents in clinical medicine documented twenty studies spanning discrete
clinical tasks including diagnostic reasoning, medication management,
and evidence retrieval, with agent systems outperforming their base
model counterparts by a median of 53 percentage points in single-agent
tool-calling studies, underscoring both the clinical promise and the
governance urgency of agentic deployment~{[}2{]}.

\subsection{Why This Changes the Governance
Problem}\label{why-this-changes-the-governance-problem}

The governance implications of this shift are not incremental. They are
structural. When AI acts as a single-inference tool, accountability is
straightforward: a clinician requests a prediction, evaluates it, and
decides whether to act on it. The AI is a consultant whose opinion is
accepted or rejected. When AI acts as an agent, this structure
dissolves. The agent is no longer a consultant. It is a participant in
care delivery. Its actions occur at speeds that preclude human review at
each step. Its errors compound: an incorrect intermediate decision at
step two affects all subsequent steps. Responsibility becomes diffuse
across the AI system, the developers, the institution, and the
clinicians nominally supervising the process.

\subsection{The Governance Gap}\label{the-governance-gap}

Existing clinical AI governance frameworks were not designed for this
environment. TRIPOD+AI~{[}3{]} provides reporting standards for model
development and validation but addresses models, not pipelines.
DECIDE-AI~{[}4{]} guides the evaluation of clinical decision-support
interventions but does not distinguish between single-inference tools
and multi-step agents. The FDA's AI/ML-based Software as a Medical
Device (SaMD) framework~{[}5{]} addresses pre-market validation and
post-market surveillance but was conceived before agentic deployment
became clinically common. The EU AI Act~{[}6{]} imposes requirements for
high-risk AI systems in healthcare but provides limited operational
guidance for autonomous pipeline governance. None of these frameworks
specifies how to measure, record, or verify trustworthiness at the level
of individual agent decisions in a running clinical system.

This paper addresses that gap. We propose a five-layer trustworthiness
governance framework specifically designed for agentic healthcare AI,
and argue that distributed ledger technologies provide the
accountability infrastructure necessary to make this framework
operational and verifiable at scale.


\section{What Makes Agentic Healthcare AI
Different}\label{what-makes-agentic-healthcare-ai-different}

Before governance frameworks can be designed, the distinctive properties
of agentic AI that create governance challenges must be understood. Four
properties differentiate agentic AI from traditional clinical prediction
models.

\subsection{Sequential Decision-Making and Error
Compounding}\label{sequential-decision-making-and-error-compounding}

Traditional clinical AI makes a single prediction given a single input.
Agentic AI makes sequences of decisions, where each step's output feeds
the next step's input. This creates an error compounding dynamic absent
from single-inference systems: an overconfident probability estimate at
step two, if not flagged and corrected, propagates into step four as if
it were reliable. Calibration, the alignment between predicted
probability and true event frequency, matters not only at the final
output but at every intermediate decision that affects subsequent steps.

Current calibration evaluation methodology was developed for
single-output models. A model's Expected Calibration Error (ECE) is
computed over its test set predictions as a whole~{[}7{]}. In an agentic
pipeline, each decision node has its own calibration profile, and
aggregate pipeline reliability depends on the product of individual node
reliabilities. A pipeline composed of five individually well-calibrated
decisions still produces systematically miscalibrated outputs if
calibration error accrues in a consistent direction across nodes~{[}8{]}.

\subsection{Opacity of Reasoning
Chains}\label{opacity-of-reasoning-chains}

Single-inference models are analyzed using well-established
interpretability methods: SHAP values decompose a prediction into
feature contributions, calibration curves reveal probability
reliability, and permutation importance identifies the features driving
a classification. Agentic reasoning chains, particularly those involving
large language model (LLM) components, iterative retrieval, or
multi-agent orchestration, are substantially harder to audit. The
intermediate states between initial input and final recommendation are
partially or fully inaccessible, and their contribution to the final
output often cannot be recovered after the fact.

This opacity creates an accountability deficit. When an agentic system
produces a harmful recommendation, investigators are often unable to
identify which reasoning step introduced the error, whether the error
reflected a miscalibrated prediction, a biased training distribution, or
an inappropriate action at a pipeline junction. Without a durable,
auditable record of intermediate states, post-incident analysis becomes
speculative.

\subsection{Speed Beyond Human Review
Cycles}\label{speed-beyond-human-review-cycles}

A defining property of agentic AI is its capacity to act at speeds
beyond human review. A triage agent processing 50 patient deterioration
signals per hour cannot be meaningfully supervised by a clinician
reviewing each action in real time. Meaningful human oversight requires
mechanisms beyond real-time review: governance must be embedded in the
system's design through pre-deployment constraints, hardcoded
boundaries, and automatic escalation triggers, then verified through
post-hoc audit rather than concurrent human review~{[}9{]}.

\subsection{Diffuse Moral
Responsibility}\label{diffuse-moral-responsibility}

When a single prediction model is used by a clinician and a patient is
harmed, responsibility analysis proceeds along a recognizable chain: the
model developer, the deploying institution, and the clinician who acted
on the recommendation. Agentic pipelines involving multiple AI
components, orchestration layers, external data sources, and automated
actions dissolve this chain. Responsibility becomes distributed across
parties who each have partial control over some steps and no control
over others~{[}10{]}. This diffusion is not merely a legal problem. It
creates practical barriers to learning from failures, holding actors
accountable, and deterring unsafe deployment.


\section{The Trustworthiness Gap in Current Governance
Frameworks}\label{the-trustworthiness-gap-in-current-governance-frameworks}

\subsection{What Existing Frameworks
Address}\label{what-existing-frameworks-address}

Several frameworks have emerged to govern clinical AI quality. TRIPOD+AI~{[}3{]}
extends the TRIPOD reporting guideline to include AI-specific
transparency requirements: model architecture, training data provenance,
calibration reporting, and external validation. DECIDE-AI~{[}4{]}
provides guidance for the evaluation of AI-based clinical
decision-support interventions, emphasizing clinical utility alongside
statistical performance. The FDA SaMD framework~{[}5{]} establishes a
risk-based regulatory pathway for AI-enabled medical devices, including
requirements for pre-market validation and post-market surveillance. The
EU AI Act~{[}6{]} classifies healthcare AI as high-risk and imposes
requirements for transparency, documentation, human oversight, and
robustness.

\subsection{What They Miss}\label{what-they-miss}

Despite their importance, these frameworks share a common limitation:
they were designed for model-level evaluation at a point in time, not
for ongoing pipeline-level governance of autonomous systems.

\textbf{Calibration is underspecified.} TRIPOD+AI recommends calibration
reporting but provides no threshold standards. A 2024 systematic review
and meta-analysis of prediction models for incident chronic kidney
disease derived or validated in community-based electronic health
records found calibration reporting so deficient that pooled
meta-analysis was not possible. Risk of bias was high in 64\% of
included models, and no clinical utility analyses or impact studies were
identified for any model~{[}11{]}. The implication for agentic AI
governance is direct: existing frameworks cannot ensure the probability
estimates produced by agent decision nodes are reliable enough to
support downstream automated actions.

\textbf{Uncertainty is absent.} Fewer than 4\% of clinical AI studies
report any uncertainty measure alongside their predictions~{[}12{]}. For
single-inference tools, this is a quality gap. For agentic systems where
uncertain predictions propagate into subsequent automated actions, it is
a safety gap. If an agent does not know, and cannot communicate, that
its probability estimate at step three carries high uncertainty, it
cannot appropriately defer to human review or widen its decision
threshold.

\textbf{Fairness is aggregate, not pipeline-level.} Where fairness
evaluation exists in clinical AI governance, it typically measures
whether subgroup performance gaps exist at the level of a model's test
set. It does not evaluate whether bias accrues or compounds across agent
pipeline steps, or whether demographic disparities in agent
recommendations emerge from the interaction of multiple individually
unbiased components~{[}13{]}.

\textbf{Audit trails are informal and mutable.} No current governance
framework specifies what constitutes an adequate audit record for an AI
agent's actions. Server logs, model outputs, and database entries are
institution-specific, mutable, and not independently verifiable. In a
regulatory dispute or adverse event investigation, the absence of
immutable, standardized audit records makes accountability analysis
unreliable.

\textbf{Consent frameworks predate autonomy.} Existing informed consent
mechanisms in healthcare AI were designed around patient consent to data
use and model deployment. They were not built for patient consent to
autonomous AI action on their behalf at each step of a clinical process.
As agentic AI increasingly acts without direct human mediation, the
concept of meaningful consent requires rethinking~{[}13{]}.


\section{A Five-Layer Trustworthiness Governance
Framework}\label{a-five-layer-trustworthiness-governance-framework}

We propose a five-layer framework that addresses the gaps identified
above. Each layer carries measurable criteria and applies both
pre-deployment and as a continuous monitoring protocol for running
agentic systems.

\subsection{Layer 1: Reliability Evaluation}\label{layer-1-reliability-evaluation}

Reliability evaluation governs whether the probability outputs produced
by each agent decision node are trustworthy enough to support downstream
automated action.

\textbf{Calibration assessment} requires computing the Expected
Calibration Error (ECE) for each decision node using empirical frequency
analysis across probability bins~{[}7{]}. We propose ECE $\leq 0.10$ as a
minimum deployment threshold, consistent with emerging clinical AI
quality standards. The Brier Score and Brier Skill Score provide
complementary summary measures. Guo et al.~{[}7{]} demonstrate how both
metrics reveal systematic overconfidence in deep learning models.
Critically, calibration must be evaluated not only on the original
training distribution but on the distribution of inputs the deployed
agent actually encounters, including the heterogeneous,
missing-data-rich inputs characteristic of real clinical environments.

\textbf{Uncertainty quantification} using split conformal
prediction~{[}8{]} provides distribution-free coverage guarantees at each
decision node. Conformal prediction produces a prediction set rather
than a single probability estimate. This set is guaranteed to contain
the true label with probability $1-\alpha$, where $\alpha$ is the chosen
miscoverage rate. We propose a coverage guarantee of $\geq 0.90$ for
clinical agentic AI, ensuring the agent communicates not just its best
prediction but the range of outcomes it cannot rule out. When prediction
sets include multiple labels, indicating ambiguous predictions, the
agent should defer to human review rather than proceed autonomously.

\textbf{External validation} is required before deployment: calibration
metrics should be recomputed on a prospectively collected or
geographically distinct external dataset to measure calibration drift,
the degradation of probability reliability when the model encounters a
new clinical environment~{[}11{]}.

\subsection{Layer 2: Fairness Auditing}\label{layer-2-fairness-auditing}

Fairness auditing governs whether agent reliability is equitably
distributed across patient subgroups.

Calibration should be stratified by age group (e.g., below 65 vs.~65
and older), sex, race/ethnicity where available, and clinically relevant
comorbidity groups (e.g., diabetes status, chronic kidney disease
stage). We propose a maximum allowable ECE gap of 0.05 between any two
demographic subgroups as a deployment gate. Where subgroup calibration
gaps exceed this threshold, the deploying institution should either
perform targeted recalibration on underrepresented groups or restrict
deployment to populations where calibration has been demonstrated.

False negative disparity deserves particular attention in healthcare AI,
where missed diagnoses disproportionately harm already-disadvantaged
groups~{[}13{]}. An agent well-calibrated overall but systematically
underestimating risk for elderly patients or patients from racial
minority groups amplifies existing healthcare inequities rather than
reducing them.

\subsection{Layer 3: Provenance Verification}\label{layer-3-provenance-verification}

Provenance verification governs whether the lineage of an agent's models
and training data is independently verifiable.

\textbf{Model versioning} requires each deployed agent decision node to
carry a verifiable identifier, a cryptographic hash of the model weights
and architecture, recorded at deployment. Any change to the model,
including recalibration, fine-tuning, or architecture modification, must
generate a new identifier and a new deployment record.

\textbf{Training data documentation} requires that the characteristics
of training datasets, including size, demographic composition, data
source, collection period, and preprocessing steps, be documented and
verifiable. This documentation must allow an independent reviewer to
assess whether the training distribution is representative of the
deployment environment.

\textbf{Deployment history} records when each model version was
deployed, to which clinical contexts, and under what oversight
conditions. This record enables post-incident analysis: when a patient
is harmed, investigators identify exactly which model version was
running, what it was trained on, and whether its deployment conditions
matched its validation conditions.

\subsection{Layer 4: Consent Management}\label{layer-4-consent-management}

Consent management governs the scope of autonomous AI action to which a
patient has agreed.

Current consent frameworks in healthcare AI address consent to data use
and model deployment at the institutional level. Agentic AI requires
more granular consent: a patient should be able to specify whether they
consent to autonomous AI actions affecting their care plan, whether they
require human review before AI-initiated actions take effect, and
whether they retain the right to override an AI agent's recommendation
at any step.

This requires a consent architecture capable of expressing, storing, and
enforcing these preferences, and updating them as the patient's care
context changes. Static institutional consent forms signed at admission
are insufficient for an environment where AI agents act on a patient's
behalf across multiple care episodes, data systems, and clinical teams.

\subsection{Layer 5: Deployment Readiness and Ongoing Monitoring}\label{layer-5-deployment-readiness-and-ongoing-monitoring}

Deployment readiness scoring integrates the preceding four layers into a
structured pre-deployment gate and a continuous post-deployment
monitoring protocol.

Pre-deployment, each agent is assessed against a deployment readiness
checklist:

\begin{enumerate}
\def\labelenumi{\arabic{enumi}.}
\tightlist
\item
  Discrimination adequacy (AUROC $\geq 0.85$ on external validation)
\item
  Calibration adequacy (ECE $\leq 0.10$ on external validation)
\item
  Calibration stability (calibration drift $\leq 0.05$ across internal
  and external datasets)
\item
  Uncertainty coverage (conformal coverage $\geq 0.90$ on external
  validation)
\item
  Coverage stability (coverage drift $\leq 0.05$ across datasets)
\item
  Prediction interpretability (singleton rate $\geq 0.70$, the
  proportion of instances receiving an unambiguous prediction)
\item
  Subgroup calibration equity (ECE gap across subgroups $\leq 0.05$)
\item
  Transparency and reproducibility (full code, data pipeline, and model
  documentation publicly or institutionally accessible)~{[}14{]}
\end{enumerate}

Post-deployment, ongoing monitoring tracks calibration drift over time
using sequential ECE monitoring with statistical process control alerts.
A statistically significant increase in ECE above the deployment
threshold triggers a mandatory review cycle and, if unresolved,
mandatory decommissioning.


\section{Distributed Ledger Technologies as Accountability
Infrastructure}\label{distributed-ledger-technologies-as-accountability-infrastructure}

\subsection{Why Traditional Logging Is
Insufficient}\label{why-traditional-logging-is-insufficient}

The trustworthiness framework described above generates a substantial
body of evidence: calibration metrics, fairness audits, provenance
records, consent documents, deployment readiness scores, and ongoing
monitoring data. For this evidence to support accountability in
regulatory review, adverse event investigation, or litigation, it must
be durable, verifiable, and independent of any single institution's
control.

Traditional server logs and database records fail on all three criteria.
They are mutable: records can be altered by administrators with access.
They are siloed: records from different components of a
multi-institutional agentic pipeline are stored in incompatible systems
without a common reference. They are not independently verifiable: an
institution claiming its AI agent behaved correctly at a given time
cannot prove its own internal records have not been modified.

These limitations are not merely theoretical. In healthcare environments
where AI agents act across institutional boundaries, for example a
diagnostic agent querying a regional radiology system or a triage agent
integrating data from multiple hospitals, no single institution's
logging architecture provides an authoritative audit record of the full
pipeline.

\subsection{Blockchain-Based Audit
Trails}\label{blockchain-based-audit-trails}

Blockchain and distributed ledger technologies address these limitations
through properties architecturally suited to the accountability
requirements of agentic healthcare AI. As blockchain and AI converge as
complementary pillars of digital health transformation, their
integration for governance and accountability has emerged as a priority
challenge for healthcare systems globally~{[}15{]}.

\textbf{Immutability} means that once a transaction is written to the
ledger, it cannot be altered without consensus from the network, making
post-hoc modification of audit records computationally and
organizationally prohibitive. For agentic AI governance, each
significant agent decision (a risk score above a clinical threshold, an
automated order, a care plan modification) is recorded as an on-chain
transaction, creating a tamper-resistant record of what the agent did,
when, on what inputs, and with what output.

\textbf{Verifiability} means any authorized party independently confirms
that a record has not been modified since it was written. A regulatory
auditor, a patient's legal representative, or a clinical ethics
committee verifies the audit trail of an AI agent's actions without
depending on the deploying institution's attestation.

\textbf{Decentralization} means the audit record is not controlled by
any single party. In a multi-institutional agentic pipeline, a shared
distributed ledger allows each participant to contribute to and verify
the record without any single party having unilateral control over what
the record says~{[}16{]}.

\textbf{Smart contracts} enable the encoding of governance rules as
executable code running on the ledger. Rather than relying on
institutional policy, a smart contract enforces governance rules
automatically: if an agent's conformal prediction set contains multiple
labels, the contract blocks autonomous action and generates a human
escalation event recorded on-chain. Ullah et al.~demonstrate this
approach in the EHR domain, implementing purpose-based access control
policies through smart contracts that establish comprehensive, immutable
audit trails recording transaction details, access purposes, digital
signatures, and timestamps, providing accountability for all data access
events while eliminating dependence on trusted intermediaries~{[}16{]}.
Governance rules become operational constraints with independently
verifiable enforcement, not aspirational policies.

\subsection{Decentralized Consent
Management}\label{decentralized-consent-management}

Patient consent for autonomous AI action is particularly well-suited to
blockchain-based management. A consent record written to a distributed
ledger is patient-controlled, institution-independent, and auditable. A
patient who consents to AI-assisted triage but not AI-initiated
medication changes records this granular preference in a consent entry
that any component of the agentic pipeline checks before acting.

Smart contracts enforce consent boundaries in real time: when an AI
agent proposes an action, the contract checks whether that action class
falls within the patient's recorded consent scope and blocks it if not,
generating a logged exception for human review. This architecture gives
consent teeth. Consent is not simply a document signed at admission but
a live constraint on system behavior.

\subsection{Verifiable Model
Provenance}\label{verifiable-model-provenance}

The provenance requirements described in Layer~3 are directly
implementable using on-chain model registries. When a new model version
is deployed, its cryptographic hash, computed from the model weights,
architecture specification, and training data characteristics, is
written to the ledger. Any subsequent query to the agent verifies, in
real time, that the model responding to the query is the registered,
validated version and not a modified derivative.

This is particularly important in environments where models are
periodically updated: a post-deployment recalibration that improves
performance on a new patient population sometimes degrades performance
on a previously studied subgroup. An on-chain model registry makes every
version change traceable, reversible, and independently verifiable. Liu
and Hu's systematic review of blockchain-AI integration in healthcare,
spanning 1,221 publications from 2015 to 2025, identifies model
provenance, traceability, and smart contract-enforced data governance as
foundational requirements for building a trustworthy and intelligent
healthcare ecosystem, with blockchain providing the immutability and
auditability that AI-generated clinical outputs require~{[}17{]}.

\subsection{Federated Learning for Privacy-Preserving Distributed
Governance}\label{federated-learning-for-privacy-preserving-distributed-governance}

Governance at scale, across multiple hospitals, health systems, and
jurisdictions, requires access to diverse patient populations for
calibration validation and fairness auditing. Federated learning enables
models to be trained and calibrated across distributed datasets without
centralizing sensitive patient data, with distributed ledger
technologies providing the coordination, verification, and audit
infrastructure for the federated process~{[}18{]}.

Each participating institution trains on its local data, contributes
gradient updates or model statistics to the federated process, and
receives in return a globally calibrated model whose performance has
been verified across the federation's full population diversity. The
ledger records each institution's participation, the aggregation
protocol used, and the resulting model provenance, providing
accountability for the full federated governance process.


\section{Implementation Challenges}\label{implementation-challenges}

The framework proposed here is architecturally sound, but its
implementation faces challenges that a research agenda must address
honestly.

\subsection{Latency and Throughput}\label{latency-and-throughput}

Writing audit transactions to a distributed ledger introduces latency
that is often incompatible with real-time clinical decision support. A
triage agent writing an on-chain transaction for each decision step
faces throughput constraints that delay time-sensitive interventions.
Architectural solutions including off-chain computation with on-chain
commitment, periodic batch anchoring of audit summaries, and layer-2
scaling protocols help mitigate this, but they require careful design to
preserve the immutability properties that motivate on-chain governance
in the first place.

\subsection{Interoperability with Existing Health Information
Systems}\label{interoperability-with-existing-health-information-systems}

Healthcare institutions operate across a heterogeneous collection of
electronic health record systems, laboratory information systems, and
clinical data warehouses that predate blockchain-based architectures.
Integrating on-chain governance with existing health information
infrastructure requires FHIR-compliant data exchange, HL7 interface
standards, and institutional commitment to adopting new auditing
protocols. On-chain governance cannot be imposed on existing systems
without substantial integration engineering.

\subsection{GDPR and the Right to
Erasure}\label{gdpr-and-the-right-to-erasure}

Blockchain's immutability creates a direct tension with GDPR's right to
erasure. A patient who invokes the right to have their data deleted
cannot have that data removed from an immutable ledger. One
architectural approach stores only cryptographic commitments on-chain
while keeping personal data off-chain. Erasure is implemented by
deleting the off-chain data and rendering the commitment unresolvable,
preserving both the audit record's integrity and the practical effect of
erasure. Bai et al.~demonstrate this hybrid model in practice for smart
healthcare systems, using IPFS for off-chain storage of personally
identifiable and protected health information while maintaining only
encrypted references on the blockchain ledger, implementing GDPR Article
17 compliance by enabling deletion from IPFS storage without
compromising the integrity of the on-chain audit record~{[}19{]}. The
broader legal adequacy of these hybrid approaches across different
regulatory jurisdictions remains an open question.

\subsection{Institutional Adoption and
Governance}\label{institutional-adoption-and-governance}

Distributed ledger governance requires consortium formation among
participating institutions: agreement on ledger architecture, consensus
mechanisms, access control, and governance of the governance layer
itself. Healthcare institutions accustomed to managing their own data
infrastructure often resist participation in a shared ledger they do not
unilaterally control. Building the institutional trust and regulatory
legitimacy for cross-institutional blockchain-based AI governance is a
sociotechnical challenge that technical architecture alone cannot solve.

\subsection{Clinician Trust
Calibration}\label{clinician-trust-calibration}

The governance framework described here is designed primarily around
system-level accountability. Agentic AI governance also has a human
dimension: clinicians must be neither over-reliant on AI agent
recommendations nor systematically distrustful of them. Calibration, in
this context, refers not only to probability reliability but to the
alignment between a clinician's confidence in an AI recommendation and
the AI's actual reliability. Governance frameworks that improve system
trustworthiness must also invest in clinician education about when and
how to appropriately weight AI recommendations.


\section{Research Agenda}\label{research-agenda}

The framework proposed here is a conceptual contribution that requires
empirical validation. Five research priorities follow from it.

\textbf{Priority 1: Empirical calibration evaluation of deployed
clinical AI.} The framework's reliability layer requires that
calibration be measured and documented for each agent decision node.
Large-scale empirical studies measuring ECE, Brier scores, and conformal
coverage rates across deployed clinical AI systems, including agentic
pipelines, would establish the current state of the field and identify
where governance intervention is most urgently needed. The finding that
calibration reporting was insufficient for meta-analysis even in a 2024
systematic review of CKD prediction models~{[}11{]} illustrates how far
current practice falls from the governance standard the framework
requires.

\textbf{Priority 2: Smart contract frameworks for clinical AI ethical
constraints.} The formal specification of governance rules in smart
contract code, particularly escalation triggers, consent enforcement,
and calibration drift alerts, requires the development of
healthcare-specific smart contract templates validated against clinical
governance requirements. Piloting these frameworks in simulated clinical
environments would provide evidence of operational feasibility.

\textbf{Priority 3: Cross-institutional federated governance pilots.}
The federated learning and distributed governance architecture described
in Section~5.5 requires multi-institutional pilots to evaluate
feasibility, latency, interoperability, and governance dynamics. Such
pilots would also generate the diverse external validation datasets
needed to properly evaluate calibration stability and subgroup fairness
across institutional settings.

\textbf{Priority 4: Patient-facing uncertainty communication standards.}
If conformal prediction is to be used to communicate uncertainty to
patients and clinicians in a way that supports informed consent and
meaningful oversight, standards are needed for how uncertainty should be
displayed, interpreted, and acted upon in clinical interfaces. Human
factors research on uncertainty communication in clinical decision
support would bridge the technical and human dimensions of the
governance framework.

\textbf{Priority 5: Prospective evaluation of DLT audit trails in
adverse event investigation.} The value of blockchain-based audit trails
ultimately depends on evidence from actual adverse event investigations.
Prospective studies comparing the completeness, accuracy, and usefulness
of on-chain versus conventional audit records in clinical AI incident
reviews would provide the evidence base needed to justify the
significant investment required for DLT adoption in healthcare AI
governance.


\section{Conclusion}\label{conclusion}

Agentic AI in healthcare is not a future prospect. It is a present
reality. Systems that autonomously monitor, reason, and act on behalf of
patients are already in clinical use, and their capabilities are
expanding rapidly. The governance frameworks currently available to
healthcare institutions, regulators, and patients were designed for a
simpler world: a model produces a prediction, a clinician evaluates it,
a patient receives care. AI agents now make sequential decisions at
machine speed. Those frameworks leave critical accountability gaps
unfilled.

The five-layer trustworthiness governance framework proposed here
addresses those gaps with measurable, operationalizable criteria. It
integrates reliability evaluation, fairness auditing, provenance
verification, consent management, and deployment readiness scoring.
Distributed ledger technologies provide the accountability
infrastructure that makes this framework verifiable: immutable audit
trails, smart contract-enforced constraints, decentralized consent
registries, and on-chain model provenance records transform governance
from aspirational policy into operational architecture.

Trustworthiness in agentic healthcare AI cannot be certified once at
deployment and assumed thereafter. It must be earned continuously,
measured rigorously, and verified independently. The framework presented
here offers one path toward that goal, alongside a research agenda for
the work required to validate it.


\section*{Acknowledgments}\label{acknowledgments}

The author acknowledges the growing body of clinical AI governance
scholarship that informed this perspective, particularly the TRIPOD+AI,
DECIDE-AI, and emerging conformal prediction frameworks for clinical
uncertainty quantification. Portions of this manuscript were drafted
with the assistance of Claude (Anthropic), an AI language model, used
for structural drafting, editorial refinement, and reference formatting.
The author takes full responsibility for all content, interpretations,
and conclusions.


\section*{Financial Disclosures}\label{financial-disclosures}

\textbf{Funding:} None.

\textbf{Conflicts of interest:} The author declares no financial or
non-financial conflicts of interest relevant to the content of this
manuscript.


\section*{Contributors}\label{contributors}

Michael O. Eniolade contributed to conceptualization, framework
development, formal analysis, writing -- original draft, and writing --
review and editing.


\section*{Data Availability Statement}\label{data-availability-statement}

No new data were generated or analyzed in this study. This is a
conceptual perspective paper. All supporting evidence is drawn from
published literature cited in the references.


\section*{References}\label{references}

\begin{enumerate}
\def\labelenumi{\arabic{enumi}.}
\item
  Rajpurkar P, Chen E, Banerjee O, Topol EJ. AI in health and medicine.
  Nat Med. 2022;28(1):31--38.
  \url{https://doi.org/10.1038/s41591-021-01614-0}
\item
  Gorenshtein A, Omar M, Glicksberg BS, Nadkarni GN, Klang E. AI Agents
  in Clinical Medicine: A Systematic Review. medRxiv.
  2025;2025.08.22.25334232.
  \url{https://doi.org/10.1101/2025.08.22.25334232}
\item
  Collins GS, Moons KGM, Dhiman P, et al. TRIPOD+AI statement: updated
  guidance for reporting clinical prediction models that use regression
  or machine learning methods. BMJ. 2024;385:e078378.
  \url{https://doi.org/10.1136/bmj-2023-078378}
\item
  Vasey B, Nagendran M, Campbell B, et al. Reporting guideline for the
  early stage clinical evaluation of decision support systems driven by
  artificial intelligence: DECIDE-AI. BMJ. 2022;377:e070904.
  \url{https://doi.org/10.1136/bmj-2022-070904}
\item
  U.S.\ Food and Drug Administration. Artificial Intelligence/Machine
  Learning (AI/ML)-Based Software as a Medical Device (SaMD) Action
  Plan. Silver Spring, MD: FDA; 2021.
\item
  European Parliament and Council. Regulation (EU) 2024/1689 of the
  European Parliament and of the Council --- Artificial Intelligence
  Act. Official Journal of the European Union. 2024.
\item
  Guo C, Pleiss G, Sun Y, Weinberger KQ. On calibration of modern neural
  networks. Proc 34th Int Conf Mach Learn. 2017;70:1321--1330.
\item
  Angelopoulos AN, Bates S. A gentle introduction to conformal
  prediction and distribution-free uncertainty quantification. Found
  Trends Mach Learn. 2023;16(4):494--591.
  \url{https://doi.org/10.1561/2200000101}
\item
  Choudhury A, Asan O. Role of artificial intelligence in patient safety
  outcomes: systematic literature review. JMIR Med Inform.
  2020;8(7):e18599. \url{https://doi.org/10.2196/18599}
\item
  Morley J, Cowls J, Taddeo M, Floridi L. Ethical guidelines for AI in
  public service. Gov Inf Q. 2021;38(3):101577.
  \url{https://doi.org/10.1016/j.giq.2021.101577}
\item
  Haris M, Raveendra K, Travlos CK, Lewington A, Wu J, Shuweidhi F,
  Nadarajah R, Gale CP. Prediction of incident chronic kidney disease in
  community-based electronic health records: a systematic review and
  meta-analysis. Clin Kidney J. 2024;17(5):sfae098.
  \url{https://doi.org/10.1093/ckj/sfae098}
\item
  Campagner A, Ciucci D, Cabitza F. Modeling the unknown: on the role of
  uncertainty in machine learning for healthcare. Comput Methods
  Programs Biomed. 2025.
  \url{https://doi.org/10.1016/j.cmpb.2024.108576}
\item
  Mittelstadt B, Allo P, Taddeo M, Wachter S, Floridi L. The ethics of
  algorithms: mapping the debate. Big Data Soc.
  2016;3(2):2053951716679679.
  \url{https://doi.org/10.1177/2053951716679679}
\item
  Collins GS, Dhiman P, Ma J, et al. Evaluation of clinical prediction
  models (part 1): from development to external validation. BMJ.
  2024;384:e074819. \url{https://doi.org/10.1136/bmj-2023-074819}
\item
  Dershem M, Riley~III J, Venkataraman M, Nasr J, Hussain AS. 2026
  Healthcare Predictions: AI, Blockchain, and the Rise of Decentralized
  Innovation. BHTY. 2025;8(3).
  \url{https://doi.org/10.30953/bhty.v8.475}
\item
  Ullah F, He J, Zhu N, Wajahat A, Nazir A, Qureshi S, Pathan MS,
  Dev~S. Blockchain-enabled EHR access auditing: Enhancing healthcare
  data security. Heliyon. 2024;10(16):e34407.
  \url{https://doi.org/10.1016/j.heliyon.2024.e34407}
\item
  Liu J, Hu X. Blockchain meets AI in healthcare: a review of convergent
  technologies for digital health transformation. Front Blockchain.
  2026;9:1766092. \url{https://doi.org/10.3389/fbloc.2026.1766092}
\item
  Rieke N, Hancox J, Li W, et al. The future of digital health with
  federated learning. npj Digit Med. 2020;3(1):119.
  \url{https://doi.org/10.1038/s41746-020-00323-1}
\item
  Bai P, Kumar S, Kumar K, Kaiwartya O, Mahmud M, Lloret J. GDPR
  Compliant Data Storage and Sharing in Smart Healthcare System: A
  Blockchain-Based Solution. Electronics. 2022;11(20):3311.
  \url{https://doi.org/10.3390/electronics11203311}
\end{enumerate}

\end{document}
\typeout{get arXiv to do 4 passes: Label(s) may have changed. Rerun}
